\subsection{Cayley-Hamilton and Annihilators}

\begin{theorem} \label{thm:cayley-hamilton-operator} \hlt{Cayley--Hamilton Theorem for Operators}\\
Let $\mathcal{V}\neq\{\mathbf{0}\}$ be finite-dimensional over $\mathbb{C}$, $q_T$ be characteristic polynomial of $T\in\LL(\mathcal{V})$. Then $q_T(T)=0$.
\begin{proof}
Let $\lambda_1,\ldots,\lambda_m$ be the distinct eigenvalues and set $d_j=\dim \mathcal{G}_{\lambda_j}(T)$.\\
It suffices, by primary decomposition, to show $q_T(T)$ kills each $\mathcal{G}_{\lambda_j}(T)$.\\
On $\mathcal{G}_{\lambda_j}(T)$, the operator $T-\lambda_jI_{\mathcal{V}}$ is nilpotent, so Corollary \ref{cor:nilpotent-power-dimension-zero} gives $(T-\lambda_jI_{\mathcal{V}})^{d_j}=0$ on $\mathcal{G}_{\lambda_j}(T)$.\\
Since the factors of $q_T(T)$ commute, this factor kills the whole product on $\mathcal{G}_{\lambda_j}(T)$. Hence $q_T(T)=0$.\\
\hlt{Structural Commentary}: Check the polynomial on each primary summand.\\
\hlt{Key Insight}: One factor of $q_T(T)$ kills each summand.
\end{proof}
\end{theorem}

\begin{theorem} \label{thm:cayley-hamilton-matrix} \hlt{Cayley--Hamilton Theorem}\\
Let $A\in\operatorname{Mat}(n,n,\mathbb{C})$ have characteristic polynomial $p_A(z)=\det(zI_n-A)$. Then $p_A(A)=0$.
\begin{proof}
Choose a Schur form $A=UTU^*$ with diagonal $\lambda_1,\ldots,\lambda_n$.\\
Since $p_A(A)=Up_A(T)U^*$ and $p_A(z)=\prod_{i=1}^{n}(z-\lambda_i)$, it suffices to prove $(T-\lambda_1I_n)(T-\lambda_2I_n)\cdots(T-\lambda_nI_n)=0$.\\
Let $Q_j\defeq\prod\limits_{i=1}^{j}(T-\lambda_iI_n)$. We show that the first $j$ columns of $Q_j$ are zero. For $j=1$, the first column of $T-\lambda_1I_n$ is zero as $T$ is upper triangular with first diagonal entry $\lambda_1$. Assume the first $j$ columns of $Q_j$ are zero.\\
Since $T-\lambda_{j+1}I_n$ is upper triangular and its $(j+1,j+1)$ entry is zero, its first $j+1$ columns lie in $\operatorname{span}(\mathbf{e}_1,\ldots,\mathbf{e}_j)$. By the induction hypothesis, $Q_j$ kills this span, so $Q_{j+1}=Q_j(T-\lambda_{j+1}I_n)$ has first $j+1$ zero columns.\\
At $j=n$, all columns of $Q_n$ are zero, so $Q_n=0$.\\
\hlt{Structural Commentary}: Schur form reduces the theorem to a triangular column-killing induction.\\
\hlt{Key Insight}: Each factor $T-\lambda_iI_n$ kills the next leading column.
\end{proof}
\end{theorem}

\begin{figure}[H]
  \centering
  \begin{tikzpicture}[
  >=Stealth,
  every node/.style={font=\small},
  mat/.style={draw=gray!55, fill=gray!5, rounded corners=2pt, inner sep=6pt, align=center},
  phasebox/.style={draw=blue!60!black, fill=blue!7, rounded corners=2pt, inner sep=5pt, align=center},
  arr/.style={->, line width=0.75pt, blue!70!black}
]
  \node[mat] (T) at (0,0)
    {$T=\begin{bmatrix}
      \lambda_1&*&*\\
      0&\lambda_2&*\\
      0&0&\lambda_3
    \end{bmatrix}$};
  \node[phasebox] (p1) at (3.7,0.8) {$P_1=T-\lambda_1I$\\first column $0$};
  \node[phasebox] (p2) at (3.7,-0.8) {$P_2=P_1(T-\lambda_2I)$\\first two columns $0$};
  \node[mat] (pn) at (7.7,0) {$P_n=\prod_i(T-\lambda_iI)=0$};

  \draw[arr] (T.east) -- (p1.west);
  \draw[arr] (p1.south) -- (p2.north);
  \draw[arr] (p2.east) -- (pn.west);
\end{tikzpicture}

\caption{In Schur form, the factors $T-\lambda_iI$ kill the leading columns one layer at a time.}
  \label{fig:cayley-hamilton-triangular-killing}
\end{figure}

\begin{theorem} \label{thm:minimal-polynomial-basic-properties} \hlt{Basic Properties of the Minimal Polynomial}\\
Let $A\in\operatorname{Mat}(n,n,\mathbb{C})$, and let $\lambda_1,\ldots,\lambda_d$ be the distinct eigenvalues of $A$.
\begin{enumerate}[label=(\roman*)]
\setlength{\itemsep}{0pt}
\item $\mu_A$ exists, is unique, and $\deg\mu_A\leq n$.
\item If a nonzero polynomial $f$ annihilates $A$, then $\mu_A$ divides $f$. In particular, $\mu_A$ divides $p_A$.
\item There are integers $q_i$ satisfying $1\leq q_i\leq m_A(\lambda_i)$ such that $\mu_A(z)=\prod\limits_{i=1}^{d}(z-\lambda_i)^{q_i}$.
\end{enumerate}
\begin{proof}
\begin{enumerate}[label=(\roman*)]
\setlength{\itemsep}{0pt}
\item Cayley--Hamilton supplies the monic annihilator $p_A$, so a least-degree monic annihilator $\mu_A$ exists and has degree at most $n$. If $g$ is another least-degree monic annihilator, write $g=h\mu_A+r$, where $r=0$ or $\deg r<\deg\mu_A$. Evaluating at $A$ gives $r(A)=0$. If $r\neq0$, scaling $r$ to be monic would give a smaller-degree monic annihilator, a contradiction. Hence $r=0$, so $g=h\mu_A$.\\
Since $g$ and $\mu_A$ are monic of the same degree, $h=1$ and $g=\mu_A$.
\item If $f(A)=0$, divide $f=q\mu_A+r$ with $\deg r<\deg\mu_A$. Evaluating at $A$ gives $r(A)=0$, hence $r=0$ and $\mu_A\mid f$. Taking $f=p_A$ gives $\mu_A\mid p_A$.
\item Since $\mu_A\mid p_A=\prod_{i=1}^{d}(z-\lambda_i)^{m_A(\lambda_i)}$, its roots are eigenvalues with exponents at most the algebraic multiplicities. If $A\mathbf{v}=\lambda_i\mathbf{v}$ and $\mathbf{v}\neq\mathbf{0}$, then $\mu_A(\lambda_i)\mathbf{v}=\mu_A(A)\mathbf{v}=0$, so every $\lambda_i$ occurs.
\end{enumerate}
\hlt{Structural Commentary}: Polynomial division makes $\mu_A$ the least monic generator of annihilators.\\
\hlt{Key Insight}: Cayley--Hamilton gives existence; eigenvectors force the roots.
\end{proof}
\end{theorem}

\begin{figure}[H]
  \centering
  \begin{tikzpicture}[
  >=Stealth,
  every node/.style={font=\small},
  box/.style={draw=gray!55, fill=gray!5, rounded corners=2pt, align=center, inner xsep=5pt, inner ysep=4.5pt, minimum width=1.55cm},
  main/.style={draw=blue!60!black, fill=blue!7, rounded corners=2pt, align=center, inner xsep=5pt, inner ysep=4.5pt, minimum width=1.55cm},
  arr/.style={->, line width=0.75pt, blue!70!black}
]
  \node[main] (mu) at (0,0) {$\mu_A$};
  \node[box] (char) at (2.60,0.55) {$p_A$};
  \node[box] (other) at (2.60,-0.55) {$f(A)=0$};

  \draw[arr] (mu) -- (char);
  \draw[arr] (mu) -- (other);
\end{tikzpicture}

  \caption{$\mu_A$ is the least monic annihilator; arrows indicate that $\mu_A$ divides $p_A$ and every annihilator $f$.}
  \label{fig:minimal-polynomial-divisibility}
\end{figure}

\begin{theorem} \label{thm:minimal-polynomial-direct-sum-disjoint-spectra} \hlt{Minimal Polynomial of a Spectrally Disjoint Direct Sum}\\
Let $A=A_1\oplus\cdots\oplus A_d$, where the spectra of $A_1,\ldots,A_d$ are pairwise disjoint. Then $\mu_A=\mu_{A_1}\mu_{A_2}\cdots\mu_{A_d}$.
\begin{proof}
Let $f=\prod\limits_i\mu_{A_i}$. Since $f(A_i)=0$ for all $i$, blockwise evaluation gives $f(A)=0$, and therefore $\mu_A\mid f$.\\
Conversely, $\mu_A(A_i)=0$ for each $i$, so $\mu_{A_i}\mid\mu_A$ for each $i$. The roots of $\mu_{A_i}$ are exactly $\operatorname{spec}(A_i)$ by Theorem \ref{thm:operator-minimal-polynomial-roots-eigenvalues}, so the disjointness of the spectra makes the polynomials $\mu_{A_i}$ pairwise relatively prime.\\
Hence their product $f$ divides $\mu_A$. Since $\mu_A\mid f$ and $f\mid\mu_A$, and both are monic, $\mu_A=f$.\\
\hlt{Structural Commentary}: Direct sums are annihilated blockwise; disjoint spectra make the block factors coprime.\\
\hlt{Key Insight}: Each spectral component contributes its own minimal-polynomial factor.
\end{proof}
\end{theorem}

\begin{theorem} \label{thm:similarity-preserves-minimal-polynomial} \hlt{Similarity Preserves the Minimal Polynomial}\\
If $A,B\in\operatorname{Mat}(n,n,\mathbb{C})$ are similar, then $p_A=p_B$ and $\mu_A=\mu_B$.
\begin{proof}
The identity $p_A=p_B$ is Theorem \ref{thm:similarity-preserves-characteristic-polynomial}. Let $A=SBS^{-1}$, then $f(A)=Sf(B)S^{-1}$ for every polynomial $f$, so $A$ and $B$ have the same annihilating polynomials. Hence $\mu_A=\mu_B$.\\
\hlt{Structural Commentary}: Similarity preserves polynomial equations.\\
\hlt{Key Insight}: Conjugation commutes with polynomial evaluation.
\end{proof}
\end{theorem}

\begin{theorem} \label{thm:diagonalizable-iff-minimal-polynomial-squarefree} \hlt{Diagonalizability and the Minimal Polynomial}\\
Let $A\in\operatorname{Mat}(n,n,\mathbb{C})$ have distinct eigenvalues $\lambda_1,\ldots,\lambda_d$.\\
Then $A$ is diagonalizable if and only if $\mu_A(z)=(z-\lambda_1)\cdots(z-\lambda_d)$.\\
Equivalently, $A$ is diagonalizable if and only if $\mu_A$ is square-free.
\begin{proof}
Set $p(z)\defeq\prod_{i=1}^{d}(z-\lambda_i)$.
\begin{enumerate}[label=\roman*.]
\setlength{\itemsep}{0pt}
\item[$(\Rightarrow)$] If $A=S\Lambda S^{-1}$ is diagonalizable, then $p(\Lambda)=0$, so $p(A)=0$.\\
Since $\mu_A$ contains every factor $z-\lambda_i$ and divides $p$, we have $\mu_A=p$.
\item[$(\Leftarrow)$] Suppose $\mu_A=p$. On the primary summand $\mathcal{G}_{\lambda_i}(A)$, the restriction of $A-\lambda_jI_n$ for $j\neq i$ is
$(A-\lambda_iI_n)+(\lambda_i-\lambda_j)I$, a nonzero scalar plus a nilpotent operator. By Lemma \ref{lem:scalar-plus-nilpotent-invertible}, these restricted factors are invertible. Since $0=p(A)|_{\mathcal{G}_{\lambda_i}(A)} =\left(\prod\limits_{j\neq i}(A-\lambda_jI_n)\right)|_{\mathcal{G}_{\lambda_i}(A)} (A-\lambda_iI_n)|_{\mathcal{G}_{\lambda_i}(A)}$, the invertible product of the $j\neq i$ factors can be cancelled, giving $(A-\lambda_iI_n)=0$ on $\mathcal{G}_{\lambda_i}(A)$.\\
Thus every primary summand is an eigenspace, so $A$ is diagonalizable.
\end{enumerate}
\hlt{Structural Commentary}: Repeated factors of $\mu_A$ record nilpotent parts in primary blocks.\\
\hlt{Key Insight}: Square-free $\mu_A$ kills each primary block in one step.
\end{proof}
\end{theorem}

\begin{corollary} \label{cor:squarefree-annihilator-implies-diagonalizable} \hlt{Square-Free Annihilator Criterion}\\
Let $A\in\operatorname{Mat}(n,n,\mathbb{C})$. If a polynomial $f$ with no repeated roots satisfies $f(A)=0$, then $A$ is diagonalizable.
\begin{proof}
Since $\mu_A\mid f$, the minimal polynomial is square-free; Theorem \ref{thm:diagonalizable-iff-minimal-polynomial-squarefree} applies.\\
\hlt{Structural Commentary}: A square-free annihilator forces square-free $\mu_A$.\\
\hlt{Key Insight}: Square-free $\mu_A$ means diagonalizable.
\end{proof}
\end{corollary}

\begin{method} \label{meth:operator-minimal-polynomial-from-matrix} \hlt{Computing $\mu_T$ from a Matrix}\\
Let $A\in\operatorname{Mat}(n,n,\mathbb{F})$ be the matrix of $T$ in some basis.
\begin{enumerate}[label=(\roman*)]
\setlength{\itemsep}{0pt}
\item For $m=1,2,\ldots$, test the monic candidate $p_m(z)=a_0+a_1z+\cdots+a_{m-1}z^{m-1}+z^m$.
\item Solve linear system by comparing matrix entries in $a_0I+a_1A+\cdots+a_{m-1}A^{m-1}=-A^m$ for $a_0,\ldots,a_{m-1}$.
\item If no solution exists, increase $m$. The first solvable candidate is $\mu_T$.
\end{enumerate}
\end{method}

\begin{figure}[H]
  \centering
  \begin{tikzpicture}[
  >=Stealth,
  every node/.style={font=\small},
  box/.style={draw=gray!60, fill=gray!6, rounded corners=2pt, align=center, inner xsep=5pt, inner ysep=4.5pt, minimum height=0.74cm, text width=1.9cm},
  main/.style={draw=blue!60!black, fill=blue!6, rounded corners=2pt, align=center, inner xsep=5pt, inner ysep=4.5pt, minimum height=0.74cm, text width=1.9cm},
  good/.style={draw=dartmouthgreen!70!black, fill=dartmouthgreen!8, rounded corners=2pt, align=center, inner xsep=5pt, inner ysep=4.5pt, minimum height=0.74cm, text width=1.5cm},
  arr/.style={->, line width=0.7pt, draw=black!70}
]
  \node[box, text width=2.05cm] (candidate) at (0,0)
    {monic candidate\\$p_m(z)$};
  \node[main] (system) at (2.82,0)
    {solve for\\$a_0,\ldots,a_{m-1}$};
  \node[box] (first) at (5.63,0)
    {first solvable\\degree $m$};
  \node[good] (mu) at (8.00,0)
    {$\mu_T=p_m$};

  \draw[arr] (candidate) -- (system);
  \draw[arr] (system) -- (first);
  \draw[arr] (first) -- (mu);
\end{tikzpicture}

  \caption{Solve $a_0I+\cdots+a_{m-1}A^{m-1}=-A^m$; the first solvable monic candidate is $\mu_T$.}
  \label{fig:minimal-polynomial-linear-system}
\end{figure}

\begin{method} \label{meth:minimal-polynomial-testing-from-characteristic-polynomial} \hlt{Testing a Minimal Polynomial from $p_A$}\\
Suppose $p_A(z)=\prod_i(z-\lambda_i)^{a_i}$ is known.
\begin{enumerate}[label=\roman*.]
\setlength{\itemsep}{0pt}
\item Write each candidate as $m_{\mathbf{q}}(z)=\prod_i(z-\lambda_i)^{q_i}$, where $1\leq q_i\leq a_i$.
\item Test exponent vectors $\mathbf{q}=(q_i)$ in increasing total degree $\sum_iq_i$.
\item For each $\mathbf{q}$, compute $m_{\mathbf{q}}(A)=\prod_i(A-\lambda_iI)^{q_i}$.
\item The first candidate with $m_{\mathbf{q}}(A)=0$ is $\mu_A$, by Theorem \ref{thm:minimal-polynomial-basic-properties}.
\end{enumerate}
\end{method}

\begin{theorem} \label{thm:operator-annihilator-iff-multiple-minimal-polynomial} \hlt{Annihilators Are Multiples of $\mu_T$}\\
Let $\mathcal{V}\neq\{\mathbf{0}\}$ be finite-dim, $T\in\LL(\mathcal{V})$, and $f\in\mathbb{F}[z]$ be nonzero. Then $f(T)=0$ if and only if $\mu_T$ divides $f$.
\begin{proof}
\begin{enumerate}[label=\roman*.]
\setlength{\itemsep}{0pt}
\item[$(\Rightarrow)$] Suppose $f(T)=0$. Divide $f=s\mu_T+r$, where $r=0$ or $\deg r<\deg\mu_T$.\\
Then $r(T)=0$, so minimality forces $r=0$. Hence $\mu_T$ divides $f$.
\item[$(\Leftarrow)$] If $f=\mu_Ts$, then $f(T)=\mu_T(T)s(T)=0$.
\end{enumerate}
\hlt{Structural Commentary}: Polynomial division reduces every annihilator to its remainder modulo $\mu_T$.\\
\hlt{Key Insight}: Minimality forces that remainder to be zero.
\end{proof}
\end{theorem}

\begin{corollary} \label{cor:operator-characteristic-multiple-minimal-polynomial} \hlt{$q_T$ Is a Multiple of $\mu_T$}\\
Let $\mathcal{V}\neq\{\mathbf{0}\}$ be finite-dimensional over $\mathbb{C}$ and let $T\in\LL(\mathcal{V})$.\\
Then the characteristic polynomial $q_T$ is a polynomial multiple of the minimal polynomial $\mu_T$.
\begin{proof}
Theorem \ref{thm:cayley-hamilton-operator} gives $q_T(T)=0$. Theorem \ref{thm:operator-annihilator-iff-multiple-minimal-polynomial} then implies that $\mu_T$ divides $q_T$.\\
\hlt{Structural Commentary}: Cayley--Hamilton turns the characteristic polynomial into an annihilator.\\
\hlt{Key Insight}: The minimal polynomial divides every annihilator, hence divides $q_T$.
\end{proof}
\end{corollary}

\begin{theorem} \label{thm:operator-minimal-polynomial-roots-eigenvalues} \hlt{Roots of $\mu_T$ Are Eigenvalues}\\
Let $\mathcal{V}\neq\{\mathbf{0}\}$ be finite-dimensional and let $T\in\LL(\mathcal{V})$. Then the roots of $\mu_T$ are precisely the eigenvalues of $T$.
\begin{proof}
\begin{enumerate}[label=\roman*.]
\setlength{\itemsep}{0pt}
\item[$(\Rightarrow)$] If $\mu_T(\lambda)=0$, write $\mu_T(z)=(z-\lambda)g(z)$. Since $\deg g<\deg\mu_T$, minimality gives $g(T)\neq0$; choose $\mathbf{v}$ with $\mathbf{u}=g(T)\mathbf{v}\neq\mathbf{0}$. Then $(T-\lambda I_{\mathcal{V}})\mathbf{u}=\mu_T(T)\mathbf{v}=\mathbf{0}$.
\item[$(\Leftarrow)$] If $T\mathbf{v}=\lambda\mathbf{v}$ with $\mathbf{v}\neq\mathbf{0}$, then $0=\mu_T(T)\mathbf{v}=\mu_T(\lambda)\mathbf{v}$, so $\mu_T(\lambda)=0$.
\end{enumerate}
\hlt{Structural Commentary}: A root splits off one shifted operator.\\
\hlt{Key Insight}: The quotient cannot already annihilate $T$.
\end{proof}
\end{theorem}

\begin{corollary} \label{cor:operator-minimal-characteristic-same-roots} \hlt{Minimal and Characteristic Polynomials Have the Same Roots}\\
Let $\mathcal{V}\neq\{\mathbf{0}\}$ be finite-dimensional over $\mathbb{C}$ and let $T\in\LL(\mathcal{V})$.\\
Then $\{\lambda\in\mathbb{C}:\mu_T(\lambda)=0\} = \{\lambda\in\mathbb{C}:q_T(\lambda)=0\} = \operatorname{spec}(T)$.
\begin{proof}
By Theorem \ref{thm:operator-minimal-polynomial-roots-eigenvalues}, $\{\lambda\in\mathbb{C}:\mu_T(\lambda)=0\}=\operatorname{spec}(T)$.\\
By Proposition \ref{prop:operator-characteristic-polynomial-degree-roots}, $\{\lambda\in\mathbb{C}:q_T(\lambda)=0\}=\operatorname{spec}(T)$.\\
\hlt{Structural Commentary}: Different multiplicities, same spectral scalars.\\
\hlt{Key Insight}: Minimal polynomial roots and characteristic polynomial roots are both just eigenvalues.
\end{proof}
\end{corollary}

\begin{perspective} \label{per:characteristic-versus-minimal-polynomial} \hlt{Characteristic vs. Minimal Polynomial}\\
The characteristic polynomial records every primary block with its dimension. The minimal polynomial records the least polynomial needed to annihilate all blocks at once.
\end{perspective}

\begin{theorem} \label{thm:annihilating-polynomials-control-spectrum} \hlt{Annihilating Polynomials Control the Spectrum}\\
Let $\mathcal{V}\neq\{\mathbf{0}\}$ be finite-dimensional, let $T\in\LL(\mathcal{V})$, and let $p$ be a nonconstant annihilating polynomial for $T$.
\begin{enumerate}[label=(\roman*)]
\setlength{\itemsep}{0pt}
\item Every eigenvalue of $T$ is a root of $p$.
\item If $\mathbb{F}=\mathbb{C}$, then at least one root of $p$ is an eigenvalue of $T$.
\end{enumerate}
\begin{proof}
\begin{enumerate}[label=(\roman*)]
\setlength{\itemsep}{0pt}
\item If $T\mathbf{v}=\lambda\mathbf{v}$ with $\mathbf{v}\neq\mathbf{0}$, then Proposition \ref{prop:polynomial-eigenvector-evaluation} gives $0=p(T)\mathbf{v}=p(\lambda)\mathbf{v}$. Since $\mathbf{v}\neq\mathbf{0}$, $p(\lambda)=0$.
\item Factor $p(z)=c(z-\lambda_1)\cdots(z-\lambda_m)$ over $\mathbb{C}$, with $c\neq0$. Then $0=p(T)=c(T-\lambda_1I_{\mathcal{V}})\cdots(T-\lambda_mI_{\mathcal{V}})$. At least one factor $T-\lambda_jI_{\mathcal{V}}$ is not injective; otherwise product would be injective. Hence $\lambda_j$ is an eigenvalue.
\end{enumerate}
\hlt{Structural Commentary}: Annihilating polynomial restricts spectrum to its roots; over $\mathbb{C}$, one root must appear.\\
\hlt{Key Insight}: Eigenvectors turn $p(T)$ into $p(\lambda)$.
\end{proof}
\end{theorem}

\begin{proposition} \label{prop:kernel-range-polynomial-in-operator-invariant} \label{prop:kernel-of-polynomial-in-operator-is-invariant} \hlt{Kernel and Range of a Polynomial in an Operator Are Invariant}\\
Let $T\in\LL(\mathcal{V})$ and $p\in\mathbb{F}[z]$. Then $\ker(p(T))$ and $\operatorname{range}(p(T))$ are invariant under $T$.
\begin{proof}
If $\mathbf{v}\in\ker(p(T))$, then $p(T)\mathbf{v}=\mathbf{0}$.\\
Since $p(T)$ commutes with $T$ by Proposition \ref{prop:operator-polynomial-products}, $p(T)(T\mathbf{v})=T(p(T)\mathbf{v})=T\mathbf{0}=\mathbf{0}$.\\
If $\mathbf{v}\in\operatorname{range}(p(T))$, write $\mathbf{v}=p(T)\mathbf{u}$. Then $T\mathbf{v}=Tp(T)\mathbf{u}=p(T)(T\mathbf{u})$, so $T\mathbf{v}\in\operatorname{range}(p(T))$.\\
\hlt{Structural Commentary}: Polynomial kernels and ranges are stable under the original operator.\\
\hlt{Key Insight}: Since $T$ commutes with $p(T)$, it preserves both $\ker p(T)$ and $\operatorname{range}p(T)$.
\end{proof}
\end{proposition}

\begin{theorem} \label{thm:sylvester-linear-matrix-equation} \hlt{Sylvester Linear Matrix Equation}\\
Let $A\in\operatorname{Mat}(m,m,\mathbb{C})$ and $B\in\operatorname{Mat}(n,n,\mathbb{C})$. If $\operatorname{spec}(A)\cap\operatorname{spec}(B)=\emptyset$, then for every $C\in\operatorname{Mat}(m,n,\mathbb{C})$ the equation $AX-XB=C$ has a unique solution $X\in\operatorname{Mat}(m,n,\mathbb{C})$.
\begin{proof}
Let $L(X)=AX-XB$. Since $L$ has equal finite-dimensional domain and codomain, it suffices to show $\ker L=\{0\}$. If $AX=XB$, then $p(A)X=Xp(B)$ for every polynomial $p$. \\
With $p=p_B$, Cayley--Hamilton gives $p_B(A)X=Xp_B(B)=0$. For every $\lambda\in\operatorname{spec}(A)$, disjointness gives $p_B(\lambda)\neq0$; hence $p_B(A)$ is invertible by Corollary \ref{cor:polynomial-in-matrix-invertible-spectrum-test}. Thus $X=0$, so $L$ is bijective.\\
\hlt{Structural Commentary}: Polynomial calculus pass through homogeneous solutions; disjoint spectra invert $p_B(A)$.\\
\hlt{Key Insight}: Spectral separation turns an intertwining relation into the zero relation.
\end{proof}
\end{theorem}

\begin{figure}[H]
  \centering
  \begin{tikzpicture}[
  >=Stealth,
  every node/.style={font=\small},
  box/.style={draw=gray!55, fill=gray!5, rounded corners=2pt, align=center, inner xsep=5pt, inner ysep=4.5pt},
  good/.style={draw=blue!60!black, fill=blue!7, rounded corners=2pt, align=center, inner xsep=5pt, inner ysep=4.5pt},
  arr/.style={->, line width=0.75pt, blue!70!black}
]
  \node[box] (gap) at (0,0) {$\operatorname{spec}(A)\cap\operatorname{spec}(B)=\emptyset$};
  \node[good] (poly) at (3.61,0) {$p_B(A)$ invertible};
  \node[box] (map) at (6.78,0) {$L^{-1}$ exists};

  \draw[arr] (gap.east) -- (poly.west);
  \draw[arr] (poly.east) -- (map.west);
\end{tikzpicture}

  \caption{Disjoint spectra make $p_B(A)$ invertible, so the Sylvester map $L(X)=AX-XB$ is invertible.}
  \label{fig:sylvester-equation-spectral-gap}
\end{figure}

\begin{corollary} \label{cor:block-upper-triangular-remove-disjoint-coupling} \hlt{Removing a Spectrally Separated Off-Diagonal Block}\\
Let $M=\begin{bmatrix}A&C\\0&B\end{bmatrix}$ and $\operatorname{spec}(A)\cap\operatorname{spec}(B)=\emptyset$. Then $M$ is similar to $A\oplus B$.
\begin{proof}
By Theorem \ref{thm:sylvester-linear-matrix-equation}, choose $X$ such that $AX-XB=-C$. Then
\begin{align}
\begin{bmatrix}I&-X\\0&I\end{bmatrix}
\begin{bmatrix}A&C\\0&B\end{bmatrix}
\begin{bmatrix}I&X\\0&I\end{bmatrix}
=
\begin{bmatrix}A&AX+C-XB\\0&B\end{bmatrix}
=
\begin{bmatrix}A&0\\0&B\end{bmatrix}.
\nonumber
\end{align}
\hlt{Structural Commentary}: A block shear cancels the off-diagonal block via Sylvester's equation.\\
\hlt{Key Insight}: Spectrally separated diagonal blocks cannot carry essential upper-right coupling.
\end{proof}
\end{corollary}

\begin{theorem} \label{thm:schur-block-form-by-eigenvalue-groups} \hlt{Schur Form Grouped by Eigenvalues}\\
Let $A\in\operatorname{Mat}(n,n,\mathbb{C})$ have distinct eigenvalues $\lambda_1,\ldots,\lambda_d$ with algebraic multiplicities $n_1,\ldots,n_d$.\\
Then $A$ is unitarily similar to a block upper-triangular matrix
\begin{align}
T=
\begin{bmatrix}
T_{11}&T_{12}&\cdots&T_{1d}\\
0&T_{22}&\cdots&T_{2d}\\
\vdots&\vdots&\ddots&\vdots\\
0&0&\cdots&T_{dd}
\end{bmatrix},
\qquad
T_{ii}\in\operatorname{Mat}(n_i,n_i,\mathbb{C}),
\nonumber
\end{align}
where each $T_{ii}$ is upper triangular with all diagonal entries equal to $\lambda_i$. Moreover, $A$ is similar to $T_{11}\oplus\cdots\oplus T_{dd}$.
\begin{proof}
Order the Schur diagonal as $\lambda_1,\ldots,\lambda_1,\lambda_2,\ldots,\lambda_2,\ldots,\lambda_d,\ldots,\lambda_d$.\\
Theorem \ref{thm:matrix-schur-triangularization-ordered} gives the displayed block Schur form. Since $\operatorname{spec}(T_{11})$ is disjoint from the trailing block spectrum, Corollary \ref{cor:block-upper-triangular-remove-disjoint-coupling} splits off $T_{11}$. Iterating splits off every diagonal block.\\
\hlt{Structural Commentary}: Grouped Schur form isolates spectral blocks; Sylvester shears split disjoint groups.\\
\hlt{Key Insight}: Cross-group upper blocks are coordinate artefacts.
\end{proof}
\end{theorem}

\begin{corollary} \label{cor:commuting-with-disjoint-blocks-has-zero-off-diagonal-blocks} \hlt{Commuting with Disjoint Blocks}\\
Let $A=A_{11}\oplus\cdots\oplus A_{dd}$, where the spectra of $A_{11},\ldots,A_{dd}$ are pairwise disjoint.\\
Partition $B=[B_{ij}]$ conformally with $A$. If $AB=BA$, then $B_{ij}=0$ for all $i\neq j$.
\begin{proof}
Equating $(i,j)$ blocks in $AB=BA$ gives $A_{ii}B_{ij}-B_{ij}A_{jj}=0$.\\
For $i\neq j$, the spectra of $A_{ii}$ and $A_{jj}$ are disjoint, so Theorem \ref{thm:sylvester-linear-matrix-equation} with $C=0$ gives $B_{ij}=0$.\\
\hlt{Structural Commentary}: A commuting matrix cannot mix direct-sum blocks with disjoint spectra.\\
\hlt{Key Insight}: The off-diagonal block solves a homogeneous Sylvester equation, whose only solution is zero.
\end{proof}
\end{corollary}

Generalized eigenvectors repair the main defect of ordinary eigenvectors: an operator may not have enough eigendirections to form a basis. The replacement is to study vectors killed by powers of shifted operators.

